% =============================================================================
\section{FIR de fase linear: tipos I--IV, janelamento e Kaiser}\label{sec:fir}
% =============================================================================

\subsection*{Teoria mínima}

\paragraph{Fase linear.} Se $h[n]$ tem simetria em torno de $M/2$
\[
  h[n]=h[M-n]\ (\text{simétrica})\qquad\text{ou}\qquad
  h[n]=-h[M-n]\ (\text{anti-simétrica}),
\]
então
\[
  H(e^{j\omega})=A(\omega)\,e^{-j\omega M/2}
  \ \Rightarrow\ \tau(\omega)=\frac{M}{2}\ \text{constante} .
\]
Um filtro FIR com fase linear é o único jeito de ter \textbf{atraso inteiro ou
meio-inteiro constante} sem distorcer o envelope. Filtros IIR só conseguem
fase aproximadamente linear (usando passa-tudo em cascata).

\begin{center}
\begin{tabular}{clll}
\toprule
Tipo & $M$ & Simetria & Restrições obrigatórias\\
\midrule
I   & par    & $h[n]=h[M-n]$    & nenhuma (serve para tudo)\\
II  & ímpar & $h[n]=h[M-n]$    & $H(\pi)=0$ (não serve p/ passa-altas)\\
III & par    & $h[n]=-h[M-n]$   & $H(0)=H(\pi)=0$ (não serve p/ passa-baixas nem altas)\\
IV  & ímpar & $h[n]=-h[M-n]$   & $H(0)=0$ (não serve p/ passa-baixas)\\
\bottomrule
\end{tabular}
\end{center}

\textbf{Por quê:} para o Tipo II, emparelhando $n$ com $M-n$ (paridade oposta,
pois $M$ é ímpar),
\[
  H(\pi)=\sum_{n=0}^{M}h[n](-1)^{n}=0
  \quad\text{exatamente, pois }(-1)^{n}+(-1)^{M-n}=(-1)^n(1+(-1)^M)=0 .
\]

\paragraph{Janelamento e o fenômeno de Gibbs.} O filtro ideal tem
$h_d[n]=\frac{\operatorname{sen}(\omega_c(n-M/2))}{\pi(n-M/2)}$
(infinito e não-causal). Truncar equivale a multiplicar por uma janela
retangular, o que produz oscilações na banda de passagem e limita a atenuação
da banda de rejeição a $\approx-21$ dB (independente de $M$ — é o lóbulo
lateral da janela, não a ordem, que manda). Janelas suaves (Hamming, Hann,
Kaiser) trocam lóbulo principal mais largo por lóbulos laterais muito
menores.

\paragraph{Janela de Kaiser (a única ``sintonizável'').}
\[
  w[n]=\frac{I_0\!\left[\beta\sqrt{1-\left(\frac{2n}{M}-1\right)^{2}}\right]}
  {I_0(\beta)},\qquad
  \beta=\begin{cases}
    0{,}1102(A-8{,}7), & A>50\ \text{dB}\\
    0{,}5842(A-21)^{0{,}4}+0{,}07886(A-21), & 21\le A\le50\\
    0, & A<21
  \end{cases}
\]
e a ordem necessária é
\[
  \boxed{\ M\approx\frac{A-8}{2{,}285\,\Delta\omega}\ }
\]
com $\Delta\omega$ a largura de transição em \textbf{rad/amostra}
($\Delta\omega=2\pi\Delta f/F_s$). O atraso é sempre $M/2$ amostras.

\subsection*{Exercícios resolvidos}

% ---------------------------------------------------------------- 13.1
\exer{Projete um passa-baixas FIR de fase linear com $\omega_c=0{,}3\pi$,
$M=20$ (21 coeficientes) usando janela de Hamming. Dê os coeficientes
principais e o atraso.}

\textbf{Solução.} $h_d[n]=\frac{\operatorname{sen}(0{,}3\pi(n-10))}{\pi(n-10)}$,
$h_d[10]=0{,}3$; multiplicando pela Hamming:
\[
  h[n]=h_d[n]\cdot\left[0{,}54-0{,}46\cos\!\Bigl(\frac{2\pi n}{20}\Bigr)\right].
\]
Primeiros coeficientes (calculados):
\[
  h[0..5]=\{\,0{,}0;\ 0{,}00293325;\ 0{,}00635175;\ 0{,}00378865;\
  -0{,}01240622;\ -0{,}03437747\,\},
\]
$h[10]=0{,}30000000$, e a soma (ganho DC) é
$\sum h[n]=1{,}00148267$. Simetria: $h[n]=h[20-n]$.
Atraso de grupo: $\tau=M/2=10$ amostras (fase linear exata). Ganho em
$\omega=\pi$ (soma alternada): $-1{,}06\times10^{-3}\approx0$ — praticamente o
esperado para um bom passa-baixas.

\resp{21 coeficientes simétricos com $h[10]=0{,}3$ e ganho DC $\approx1$;
atraso constante de $10$ amostras.}

% ---------------------------------------------------------------- 13.2
\exer{\textbf{(Desafio 05)} Com $M=21$ (Tipo II) e a mesma especificação do
exercício anterior, mostre numericamente que $H(\pi)=0$. Explique por que esse
filtro \emph{não} pode ser usado como passa-altas.}

\textbf{Solução.} Neste caso $M$ é \textbf{ímpar} e não existe amostra no
ponto de simetria ($M/2=10{,}5$):
\[
  h[10]=h[11]=0{,}28753440,\qquad h[n]=h[21-n] .
\]
A resposta em $\omega=\pi$ é a soma alternada:
\[
  H(\pi)=\sum_{n=0}^{21}h[n](-1)^{n}=0 \quad\text{(exatamente)},
\]
porque em cada par $(n,\,21-n)$ as paridades são opostas e os valores iguais,
logo $h[n](-1)^n+h[21-n](-1)^{21-n}=h[n](-1)^n(1-1)=0$.
Um passa-altas \emph{precisa} de ganho alto em $\omega=\pi$; como o Tipo II
tem um \textbf{zero obrigatório} em $\omega=\pi$, seria necessário um Tipo I
($M$ par) ou um Tipo IV com transformação adequada.

\resp{$H(\pi)=0$ exato para o Tipo II (zero forçado em $z=-1$); para
passa-altas use Tipo I ou III (com $M$ par, por exemplo $M=32$).}

% ---------------------------------------------------------------- 13.3
\exer{Projete por Kaiser um passa-baixas com $F_s=1000$ Hz, banda passante até
$100$ Hz, banda de rejeição a partir de $200$ Hz e atenuação $A=50$ dB:
calcule $\beta$, a ordem $M$ e o atraso.}

\textbf{Solução.} Largura de transição:
\[
  \Delta f=200-100=100\ \text{Hz}\ \Longrightarrow\
  \Delta\omega=2\pi\frac{100}{1000}=0{,}6283185307\ \text{rad/amostra}.
\]
Como $A=50$ dB ($>50$), usa-se $\beta=0{,}1102(A-8{,}7)$:
\[
  \beta=0{,}1102\,(50-8{,}7)=4{,}551260 .
\]
Ordem:
\[
  M\approx\frac{A-8}{2{,}285\Delta\omega}
   =\frac{42}{2{,}285\cdot0{,}6283185307}
   =29{,}253863\ \Longrightarrow\ \boxed{M=30}\ (31\ \text{coeficientes}).
\]
Corte em $\omega_c=0{,}35\pi$ (\emph{entre} $100$ e $200$ Hz: $f_c=175$ Hz —
qualquer ponto da faixa de transição serve, pois $\beta$ e $M$ dependem
apenas de $\Delta\omega$). Atraso: $M/2=15$ amostras; o ganho DC calculado
é $1{,}004532$ (o projeto de janela erra ligeiramente o ganho — normaliza-se
dividindo pela soma, se necessário).

\resp{$\beta=4{,}5513$, $M=30$ (31 coeficientes), $\Delta\omega=0{,}6283$,
atraso $15$ amostras, ganho DC $\approx1{,}005$.}
